\section{Conclusions}
\label{sec:conclusions}

We introduced automated code clone \emph{generation} as the constructive inverse of clone detection, and we presented \sys, an agentic system that synthesizes a behaviorally faithful but structurally divergent clone of a given program. \sys decomposes the task across four cooperating LLM-driven agents, library substitution, glue synthesis, first-party logic mutation, and cross-cutting verification, with a fifth agent standing in for the first two when the requirement is to remove a dependency rather than substitute one. Correctness is enforced at two tiers: a local execute-and-repair loop that guards every transformation and, for the agents that rewrite first-party logic, checks each candidate against recorded golden vectors; and a global acceptance gate that independently confirms the finished clone is a genuine migration before it can reach the next stage. The key insight is that generation, unlike detection, permits constructive verification: an LLM need not be correct in one shot, because a disciplined loop can execute, test, and repair its output until equivalence holds.

On software that estimates factory measurements and controls industrial infrastructure, this design works. Full \sys migrates 32 of 36 runs (88.9\%), against 18 (50.0\%) for a plain single-shot baseline and 20 (55.6\%) with one evidence-fed repair call added, and against 27 (75.0\%) for an ablation with the verification tier disabled. It also confirms equivalence on every one of its 36 runs, the only configuration we tested that does; the ablation reaches 31, and the two single-shot baselines 22 and 23. Verification's contribution is clearest here, in \emph{how much} \sys successfully migrates, not in how cleanly each individual acceptance turns out: an independent blind review of all 144 generated clones finds full \sys with both the most clones rated clean and the fewest rated major-or-broken of any configuration in absolute terms, though at our corpus size that is a direction worth noting rather than a settled conclusion. That gain is not free: full \sys spends roughly seventeen times the dollar cost, and thirty-three times the tokens, of the plain baseline to achieve it.

The blind review also draws a hard boundary around what our deterministic checks can promise. Every clone in this study, across all four configurations, compiles and passes its own test suite before migration; the checks that actually separate a correct migration from an incorrect one, our equivalence oracle and, especially, a fifth layer that executes the substituted factory path directly against the real library, exist precisely because the subject's own dependency-injected tests never exercise that path. Even so, the blind reviewer still finds real defects among clones that pass every deterministic check we run, full \sys included. Test-suite passage, and even our own oracle, are necessary signals for clone generation, not sufficient ones.

Three directions follow from what we observed while building this system. First, the fifth oracle layer we added mid-study caught defects that the source-level checks ahead of it, and the subject's own tests, all missed; it fires on a minority of subjects, and extending equivalence confirmation to the binary level, disassembling the compiled original and clone and reasoning over execution traces in the manner of our prior binary clone detector, is a direct next step we have not yet built. Second, the residual failure mode we see most often is concentrated at the seam between agents, an adapter generated correctly and then incompletely wired to the call sites it was written for, which is a program-analysis problem more than a generation problem, and suggests pairing each agent with a static check on the contract it is expected to honor. Third, the fifth agent we added for programs with no library to substitute earns a real behavioral check of its own, but the retries it receives when our verification tier rejects its output carry none of the concrete evidence that the library-substitution and glue-synthesis agents' retries do; closing that asymmetry is a direct, low-risk improvement to the architecture as it stands today.

A further direction follows from the corpus itself, not from the system. Python here is a proxy for the measurement-and-control logic real industrial systems run, not the language those systems actually use: real controllers are programmed in the IEC~61131-3 languages, Structured Text and Ladder Logic chief among them, and retargeting \sys's agents at those languages, evaluated against real controller toolchains rather than a Python stand-in, is the natural next step toward the domain this work is actually motivated by. A second extension follows from how equivalence is checked here: golden vectors and invariants confirm behavior at the source level, not against the physical process a controller drives, and running generated and migrated programs inside an industrial control system simulator that models plant dynamics, not just recorded I/O, would test equivalence at the level a real deployment needs. Third, this work's motivation is partly defensive: diversified code is meant to deny an attacker a uniform target. That claim is asserted here, not tested. A direct next step is studying, under simulation, whether an exploit effective against the original program still succeeds against \sys's structurally divergent clone, and separately, whether the agentic pipeline itself, its requirements input, its library-selection evidence, its repair loop, is a new attack surface a determined adversary could use to steer \sys toward a compromised substitution rather than a safer one. A fourth extension is scale: twelve subjects support the headline comparison this paper leans on, migration success, but leave most other comparisons too close to call, as Section~\ref{sec:limitations} notes; a larger corpus, whether authored the same way or drawn from real controller codebases once the language extension above is in place, is what would let those directional findings become established ones.

Finally, because the capacity to produce equivalent-but-divergent code is inherently dual-use, defensively valuable for diversification, yet potentially useful for evading similarity-based detection, \sys is designed to emit a machine-checkable record of every substitution and mutation it applies, so that each transformation remains attributable and auditable wherever this capability is deployed. Recognizing and generating clones are two faces of the same problem; coupling them, as \sys does, turns program equivalence from a question we merely answer into a property we can deliberately engineer.
